\scanpage{76}
\inserted{On} peut aussi expliciter dans notre contexte actuel l'hom
\[
f':\mathfrak g\longrightarrow\widetilde{\mathfrak g}\tag*{(53)}
\]
(s'identifiant à $\Lambda^2f$) tel que
\[
\left\{\begin{aligned}
ff'&=-2\operatorname{id}_{\mathfrak g}\\
f'f&=-2\operatorname{id}_{\widetilde{\mathfrak g}}
\end{aligned}\right.\tag*{(54)}
\]
en posant
\[
f'(\bar u)=(\operatorname{Tr}u)1-2u\tag*{(55)}
\]
où $u$ est une section loc de $M$, et $\bar u$ son image par $M\to\mathfrak g$ ; comme $\operatorname{Tr}1=2$ cette application \deleted{\unsure{en $u$ est}} $M\to M$ est bien à val. dans $\widetilde{\mathfrak g}$, et nulle sur $\mathcal O_S.1$ donc se factorise par $\mathfrak g$, enfin si $u$ est dans $\unsure{\widetilde{\mathfrak g}|\mathfrak g}$ on trouve $f'f(u)=-2u$, enfin $ff'(\bar u)=\deleted{\unsure{-2f(u)}}\overline{f'(\bar u)}=-2\bar u$ OK.

Il reste à voir comment s'explicitent en termes de $M$ les objets $X$, et $P\supset\deleted{\unsure{C}}X$. Bien sûr, on aura
\[
\left\{\begin{aligned}
P&=\check{\mathbb P}(\widetilde{\mathfrak g})\simeq\mathbb P(\mathfrak g)\\
X&=\mathbb V(q)\subset P
\end{aligned}\right.\tag*{(56)}
\]
\margin{\deleted{\uline{NB} les sections de $X$ sont \unsure{en corresp biuniv|aussi les suiv.} :}}
% slanted
% archive sheet number: 75
